% !TeX root = ../main.tex
\section{Conclusion}
\label{sec:conclusion}

We presented a behavioral simulation framework for analog compute-in-memory inference that separates algorithmic cross-instance transfer from physical precision-retention constraints. By comparing ideal pure-quantization baselines with realistic phase-change memory (PCM) physics, the framework identifies deployment frontiers before committing to a specific device implementation.

Two conclusions are central. First, the tested PCM UnitCell regime exposes a non-monotonic precision--accuracy curve: 8-bit PCM achieves the highest fresh accuracy (77.60\%) with negligible drift, 4-bit PCM achieves comparable fresh accuracy (76.68\%) but pays a 4.01~pp one-day drift penalty, and 6-bit PCM enters a D2D-sensitive transition zone with lower mean accuracy (68.44\%, seed-level std 6.03\%) despite drift stability. Second, the ideal-device ablation shows why this physical frontier should not be confused with algorithmic robustness. Fixed-instance training at 4-bit collapses to 14.64\%, whereas Ensemble Hardware-Aware Training (HAT) resamples mismatch masks and recovers fresh-instance accuracy to 86.16$\pm$0.19\%.

The resulting design rule is sequential. A system designer should first avoid fixed-instance overfitting, then clear the readout-precision cliff, and only then select the memory precision that satisfies retention constraints. Ensemble HAT addresses the first step; the ADC/D2D envelope and PCM precision ladder define the hardware-side gates. The same separation should guide future studies of memory-bound inference components, including analog KV-cache architectures, where train-time robustness and deployment-time device physics must again be evaluated independently.
